\section{Introduction}
\label{sec:introduction}

Let \(A,B\in\mathbb R^{d\times m}\), with \(A\) of full column rank and
\(1\le m<d\).  Consider the affine-invariant least-deformation problem
\begin{equation}
  \min_{P\in\mathbb S_{++}^d}
  \frac12 d_{\mathrm{AI}}(P_{\rm ref},P)^2
  \quad\text{subject to}\quad
  PA=B,
  \qquad
  \det P=\det P_{\rm ref}.
  \label{eq:introduction-main-problem}
\end{equation}
After an affine-invariant congruence and a scalar normalization, it is enough
to take \(P_{\rm ref}=I\) and \(\det P=1\).  The feasibility condition is
\(A^\top B\in\mathbb S_{++}^m\).  When \(m<d\), the action \(PA=B\) leaves a
positive-dimensional family of feasible completions, and
\cref{eq:introduction-main-problem} selects the one of minimum intrinsic
deformation.  The same problem is an affine-invariant multi-secant completion
when \(A\) contains secant directions and \(B\) contains their desired inverse
actions.

This completion problem is motivated by adaptive optimization.  Moments,
curvature surrogates, Kronecker factors, and tensor preconditioners act as
hidden positive operators that turn visible gradients into parameter motion
\cite{duchi2011adagrad,kingma2014adam,martens2015kfac,
gupta2018shampoo,morwani2024shampoo,jordan2024muon,
essentialai2025muon,crawshaw2025muonvariants}.  Observing the action of such an
operator on a low-dimensional subspace does not determine its complementary
geometry.  Problem \eqref{eq:introduction-main-problem} supplies a canonical
completion and an intrinsic measure of the deformation required to realize
the observed action.

Three issues prevent the formulation alone from giving an algorithmic result.
The feasible set is a nonlinear SPD fiber whose global topology and geodesic
geometry are initially unknown; nearest-point existence does not by itself
produce a computable optimality residual; and a dense implementation would
appear to require matrix functions in the ambient dimension \(d\).  We resolve
these issues by combining a global action-bundle theorem with an explicit
residual-gradient iteration and an exact active-subspace reduction.

The matrix problem sits inside a general variational-reduction mechanism.  For
hidden state \(x\), visible action \(y=\pi(x)\), and deformation energy \(E\),
define the infimal pushforward
\[
  (\pi_!E)(y)=\inf_{\pi(x)=y}E(x).
\]
It eliminates hidden realizations and composes across nested maps.
Convex-fiber regularity selects a canonical lift, while smooth nondegeneracy
gives
\[
  D\sigma_\star=-E_{\sigma\sigma}^{-1}E_{\sigma y},
  \qquad
  D^2(\pi_!E)
  =E_{yy}-E_{y\sigma}E_{\sigma\sigma}^{-1}E_{\sigma y}.
\]
When mechanism amplitudes enter affinely before reduction, their interaction
curvature is
\[
  D^2_{uu}(\pi_!E)=-G^*H^{-1}G\preceq0,
\]
and its mixed entries integrate to finite contrasts under the specified
intervention protocol.  In the SPD action problem, the same vertical Hessian
governs both this response and the susceptibility of the nearest completion.

Our principal realization is the determinant-one SPD action map
\[
  \pi_A(P)=PA.
\]
Every admissible fiber has a global analytic parameterization by a smaller
determinant-one SPD factor and is closed and totally geodesic under the
affine-invariant metric.  Projection of the reference therefore gives a unique
analytic canonical controller.  A strongly-convex fiber argument connects
this selection to computation, while the action bundle makes its logarithmic
residual explicit.  A sublevel first-return argument yields the radius
majorant \(L_0=\psi(D_0/\sqrt2)\), nonasymptotic iteration bounds from exact
inversion of the linear rate, and posterior controller and objective
certificates.  The same analysis propagates certified matrix-function errors
and radius majorants.

From the identity hidden coordinate, and more generally from an
active-compatible hidden coordinate, the solver has an exact realization on
\(\operatorname{span}(\operatorname{col}A,\operatorname{col}B)\), of dimension
\(r\le2m\).  It replaces the full spectral kernel by \(r\times r\) operations
and gives a strict dimension reduction when \(r<d\).  Arbitrary hidden
initialization retains the global convergence theorem but need not preserve
this active symmetry.  Nested shape-normalized actions then satisfy a CAT(0)
Pythagorean law and identify \(c_HH^{-1}\) at the sharp rank \(d-1\), provided
the scalar gauge \(c_H=(\det H)^{1/d}\) is known.  Finite inner solves obey an
observable residual-controlled perturbation of the same law.

\paragraph{Contributions.}
\begin{enumerate}[leftmargin=*,itemsep=0.2em]
  \item \textbf{Global geometry of the action constraint.}
  We construct a global analytic trivialization of \(P\mapsto PA\), prove that
  every determinant-one feasible fiber is closed and totally geodesic, and
  obtain a unique analytic affine-invariant nearest-completion section.

  \item \textbf{A global solver with posterior certificates.}
  The exact logarithmic residual is the whitened negative gradient of a
  strongly geodesically convex reduced objective.  Its exponential update
  converges globally and linearly from every initial hidden coordinate, with
  explicit iteration bounds and sharp posterior distance and objective-gap
  certificates.  A conditional inexact theorem propagates certified errors in
  the residual direction; update, factorization, and storage errors remain the
  responsibility of the numerical implementation.

  \item \textbf{Active complexity and finite identification.}
  From an active-compatible initializer, the exact iteration reduces to an
  active space of dimension \(r\le2m\), costs \(O(r^3)\) per step after
  \(O(dm^2+r^3)\) preprocessing, and applies the resulting controller in
  \(O(dr+r^2)\).  Nested projections obey exact and residual-controlled
  approximate CAT(0) laws and recover determinant-one inverse-Hessian shape
  after \(\lceil(d-1)/b\rceil\) independent width-\(b\) rounds, assuming the
  scalar determinant gauge is supplied.

  \item \textbf{Variational response and interaction.}
  Associative hidden-state elimination yields the canonical response and the
  Schur reduced Hessian in fixed product coordinates.  Affine mechanism
  perturbations give a
  negative-semidefinite interaction curvature whose action specialization is
  locally controlled by the same inverse vertical Hessian as completion
  sensitivity.
\end{enumerate}

Controlled synthetic experiments accompany the theorem chain.  They test the
posterior inequalities against high-accuracy reference solutions, compare the
full and active iterates, measure the dense spectral kernels as \(d\) grows,
and track the curvature radius as \(\lambda_{\min}(A^\top B)\downarrow0\).

Together, these results turn the intrinsic completion problem
\eqref{eq:introduction-main-problem} into a globally solvable reduced
optimization problem with observable termination criteria.  The global action
bundle links its geometry, sensitivity, algorithm, and finite-identification
law through one variational-reduction chain.

\paragraph{Organization.}
\Cref{sec:reduction,sec:interaction} develop variational reduction and its
response theory.  \Cref{sec:action-bundle} establishes the global action
geometry, \cref{sec:certified-solver} gives the solver, certificates, and operation
bounds, \cref{sec:numerical-validation} reports controlled numerical tests,
and \cref{sec:multisecant-recovery} proves finite recovery.  The appendices
contain complete proofs and a theorem-level comparison with the closest
literature.
